\section*{अध्याय \ref{ch:g12:comb} --- क्रमचय-संचय और गणना}

\begin{solution}{exo:g12:comb:1}
गुणन का सिद्धांत: $26^2 \times 10^3 \times 26^2 = 26^4 \times 10^3 = 456\,976\,000$।

बिना दोहराए चिह्नों के, चारों अक्षर भिन्न होने चाहिए ($26 \times 25 \times 24 \times 23$ तरीक़े, अक्षर वाले
स्थान क्रम से भरते हुए) और तीनों अंक भिन्न ($10 \times 9 \times 8$):
\[
26 \times 25 \times 24 \times 23 \times 10 \times 9 \times 8
= 358\,800 \times 720 = 258\,336\,000 .
\]
\end{solution}

\begin{solution}{exo:g12:comb:2}
$\dbinom83 = \dfrac{8 \times 7 \times 6}{3!} = 56$; $\dbinom{10}{8} = \dbinom{10}{2} = \dfrac{10 \times 9}{2} = 45$; \[
\frac{\binom n2}{\binom{n+1}2}
= \frac{n(n-1)/2}{(n+1)n/2} = \frac{n-1}{n+1}.
\]
\end{solution}

\begin{solution}{exo:g12:comb:3}
समिति चुनिए: $\binom{30}{4}$ तरीक़े। फिर उन 4 में से अध्यक्ष और कोषाध्यक्ष क्रम में
चुनिए: $4 \times 3 = 12$ तरीक़े। कुल
\[
\binom{30}{4} \times 12 = 27\,405 \times 12 = 328\,860 .
\]
\end{solution}

\begin{solution}{exo:g12:comb:4}
\[
(x+2)^5 = x^5 + 10x^4 + 40x^3 + 80x^2 + 80x + 32 ,
\]
\[
(1-x)^6 = 1 - 6x + 15x^2 - 20x^3 + 15x^4 - 6x^5 + x^6 .
\]
$(2x+3)^7$ में $x^3$ वाला पद $\binom{7}{3}(2x)^3\,3^4 = 35 \times 8 \times 81\, x^3$ है: गुणांक $22\,680$ है।
\end{solution}

\begin{solution}{exo:g12:comb:5}
\emph{1.} $\dbinom{52}{5} = 2\,598\,960$।

\emph{2.} ठीक एक इक्का: उसे चुनिए ($4$ तरीक़े) और $4$ ग़ैर-इक्कों से पूरा
कीजिए: $4 \times \binom{48}{4} = 4 \times 194\,580 = 778\,320$. कम से कम एक इक्का: \omterm{def:g10:proba:operations}{पूरक} से गिनकर $\binom{52}{5} - \binom{48}{5} = 2\,598\,960 - 1\,712\,304 = 886\,656$।

\emph{3.} तीन एक जैसे पत्तों की कोटि चुनिए ($13$), उनके रंग ($\binom43 = 4$), फिर जोड़े
की कोटि ($12$ बची हुई), उसके रंग ($\binom42 = 6$): $13 \times 4 \times 12 \times 6 = 3744$।
\end{solution}

\begin{solution}{exo:g12:comb:6}
$\mathrm{MATH}$ में 4 भिन्न अक्षर हैं: $4! = 24$ अक्षर-विन्यास।

$\mathrm{BANANA}$ में 6 अक्षर हैं: तीन $A$, दो $N$, एक $B$। पहले $A$ के स्थान
चुनिए ($\binom63$), फिर शेष में से $N$ के ($\binom32$), और $B$ को अंतिम जगह मिल जाती है:
\[
\binom{6}{3}\binom{3}{2} = 20 \times 3 = 60 .
\]
(समतुल्य रूप से $\frac{6!}{3!\,2!\,1!} = 60$।)
\end{solution}

\begin{solution}{exo:g12:comb:7}
\emph{बीजगणित से:}
\[
k\binom nk = \frac{k\,n!}{k!(n-k)!} = \frac{n!}{(k-1)!\,(n-k)!}
= n\,\frac{(n-1)!}{(k-1)!\bigl((n-1)-(k-1)\bigr)!} = n\binom{n-1}{k-1}.
\]
\emph{दोहरी गिनती से:} युग्म (के $k$ लोगों की समिति, और उसमें से अध्यक्ष)
गिनिए। या तो पहले समिति चुनिए ($\binom nk$) फिर उसका अध्यक्ष ($k$): $k\binom nk$ युग्म। या
पहले अध्यक्ष चुनिए ($n$ विकल्प) और फिर बचे हुए $n-1$ में से शेष $k-1$ सदस्य:
$n\binom{n-1}{k-1}$ युग्म।
\end{solution}

\begin{solution}{exo:g12:comb:8}
हर पथ में ठीक $m + n$ क़दम होते हैं, जिनमें $m$ पूर्व की ओर और $n$ उत्तर की ओर
हैं; और वह पूरी तरह उन क्षणों के समुच्चय से तय हो जाता है ($m+n$ में से) जिन
पर पूर्व की ओर क़दम रखा जाता है। ऐसे $\binom{m+n}{m}$ चुनाव हैं।
\end{solution}

\begin{solution}{exo:g12:comb:9}
$m + n$ लोगों के समुच्चय को $m$ के समूह $A$ और $n$ के समूह $B$ में बाँटिए।
किसी $k$-अवयवी उपसमुच्चय में $A$ के कुछ $j$ सदस्य होते हैं ($0 \leq j \leq k$) और $B$
के $k - j$; किसी नियत $j$ के लिए ऐसे $\binom mj \binom{n}{k-j}$ उपसमुच्चय हैं, और $j$ पर योग का
सिद्धांत वांडरमोंड की सर्वसमिका दे देता है।

$m = n = k$ के साथ:
\[
\binom{2n}{n} = \sum_{j=0}^n \binom nj \binom{n}{n-j}
= \sum_{j=0}^n \binom nj^{2},
\]
जहाँ सममिति $\binom{n}{n-j} = \binom nj$ का उपयोग हुआ है।
\end{solution}

\begin{solution}{exo:g12:comb:10}
\emph{\cref{exo:g12:comb:7} से:}
\[
\sum_{k=1}^n k\binom nk = \sum_{k=1}^n n\binom{n-1}{k-1}
= n\sum_{j=0}^{n-1}\binom{n-1}{j} = n\,2^{n-1},
\]
जो \cref{cor:g12:comb:sums} से मिलता है। \emph{अवकलन से:} $(1+x)^n = \sum_k \binom nk x^k$ का अवकलन $n(1+x)^{n-1} = \sum_k k \binom nk x^{k-1}$ देता है; अब $x = 1$ पर
मान निकालिए।
\end{solution}

\begin{solution}{pb:g12:comb:1}
\textbf{1.} $26^2 \times 10^3 = 676\,000$ पट्टिकाएँ। $\mathrm{BANANA}$: $6$ अक्षर, जिनमें $A$ तिगुना
और $N$ दुगुना है: $\frac{6!}{3!\,2!} = 60$ अक्षर-विन्यास।

\textbf{2.} $\binom{32}{5} = 201\,376$ हाथ; ठीक दो इक्कों वाले $\binom42 \binom{28}{3} = 6 \times 3\,276 = 19\,656$।

\textbf{3.} पथ $4$ बार $R$ और $3$ बार $U$ वाला शब्द है: $U$ के स्थान चुनिए:
$\binom73 = 35$।

\textbf{4.} $(1+x)^4 = 1 + 4x + 6x^2 + 4x^3 + x^4$। $x = 1$ पर: $\sum_k \binom nk = 2^n$; $x = -1$ पर: $\sum_k (-1)^k \binom nk = 0$ --- यानी \omterm{prop:g11:binom:pascal}{पास्कल त्रिभुज} के
पंक्ति-योग और एकांतर पंक्ति-योग।

\textbf{5.} $n$ में से अध्यक्ष सहित $k$ लोगों की समितियाँ: या तो समिति फिर
उसका अध्यक्ष चुनिए ($\binom nk \times k$), या अध्यक्ष फिर शेष सदस्य ($n \times \binom{n-1}{k-1}$): दोनों बराबर।
$k$ पर जोड़ने पर दाहिना पक्ष $n \sum_j \binom{n-1}{j} = n\,2^{n-1}$ हो जाता है।

\textbf{6.} $10$ तारों और $3$ छड़ों की पंक्ति ऑर्डर लिख देती है (पहली छड़
से पहले स्वाद 1 के गोले, इत्यादि); पंक्ति में $13$ चिह्न हैं और वह छड़ों के
स्थानों से तय हो जाती है: $\binom{13}{3} = 286$ ऑर्डर।

\textbf{7.} $12$ तारे, $2$ छड़ें: $\binom{14}{2} = 91$।

\textbf{8.} $x', y', z' \geq 0$ और $x' + y' + z' = 9$ के साथ: $\binom{11}{2} = 55$।

\textbf{9.} $i + j + k = 5$ वाला एकपदी $a^i b^j c^k$: $\binom72 = 21$।

\textbf{10.} सूत्र से: $3$ तारे, $1$ छड़: $\binom41 = 4$; सूची: $(3,0)$, $(2,1)$,
$(1,2)$, $(0,3)$: मेल।

\textbf{11.} ``एक जैसे'' वहाँ घुसा जहाँ यह घोषित किया गया कि ऑर्डर हर स्वाद
की \emph{गिनती} भर है --- तारों के कोई नाम नहीं होते। यदि गोले क्रम से खाए
जाएँ, तो $10$ भिन्न स्थानों में से हर एक स्वतंत्र रूप से कोई स्वाद चुन लेता
है: $4^{10} = 1\,048\,576$ क्रम --- यानी अलग प्रतिरूप और अलग संसार (\cref{met:g12:comb:model}: सदा पूछिए
\emph{क्रमित? भिन्न? पुनरावृत्ति की छूट?})।

\textbf{12.} $D_1 = 0$; $D_2 = 1$ (अदला-बदली); $D_3 = 2$ (दोनों $3$-चक्र); $D_4 = 9$।

\textbf{13.} अतिथि 1 को टोपी $k \neq 1$ मिलती है: $n - 1$ विकल्प। यदि अतिथि $k$ को
टोपी 1 मिले, तो बचे हुए $n - 2$ अतिथि अपनी टोपियों का विपर्यय करते हैं: $D_{n-2}$
तरीक़े। यदि अतिथि $k$ को टोपी 1 \emph{न} मिले, तो टोपी 1 को अतिथि $k$ की
वर्जित टोपी मान लीजिए: बचे हुए $n - 1$ अतिथि विपर्यय करते हैं: $D_{n-1}$ तरीक़े।
अतः $D_n = (n-1)(D_{n-1} + D_{n-2})$। जाँच: $D_4 = 3(2 + 1) = 9$; और $D_5 = 4(9 + 2) = 44$।

\textbf{14.} $3! = 6$ बँटवारों में से वे घटाइए जो कम से कम एक टोपी सही रखते हैं:
किसी दी हुई टोपी को तीन सही रखते हैं (हर एक $2!$, कुल $3 \times 2 = 6$), जिससे युग्म
अधिक गिन लिए जाते हैं ($3$ युग्म, हर एक $1!$) और उन्हें लौटाना पड़ता है,
और फिर तत्समक को दोबारा घटाना पड़ता है ($1$): $D_3 = 6 - 6 + 3 - 1 = 2$, अर्थात् $3!\left(1 - 1 + \frac12 - \frac16\right) = 2$। सामान्य
रूप से $D_n = n!\sum_{k=0}^{n} \frac{(-1)^k}{k!}$।

\textbf{15.} $\frac{D_5}{120} = \frac{44}{120} \approx
0.3667$, जो $\frac1\eu \approx 0.3679$ के पहले से ही पास है: एकांतर योग $1 - 1 + \frac{1}{2!} - \frac{1}{3!} + \dots$ $\eu^{-1}$ की
ओर बढ़ता है। किसी बड़ी दावत की टोपियाँ लगभग $36.8\,\%$ बार पूरी तरह बिखर जाती हैं
--- लॉटरी और सचिव वाला वही अचर, तीसरी बार दिखाई देते हुए।

\textbf{16.} $\P(\text{वैध}) = \frac{D_{10}}{10!} \approx
0.368$। हर बार फिर से निकालने पर सफलता की \omterm{def:g10:proba:distribution}{प्रायिकता} $\approx \frac1\eu$ है,
इसलिए निकालों की प्रत्याशित संख्या लगभग $\eu \approx 2.7$ है: तीन बार टोपियाँ बाँटने का
बजट रखिए।

\textbf{17.} हर हस्तमिलन कुल कोटि-गिनती में $2$ जोड़ता है, इसलिए सभी अतिथियों
की हस्तमिलन-संख्याओं का योग सम होता है। \omterm{def:g10:numbers:sets}{पूर्णांकों} का योग तभी सम होता है जब
विषम पदों की संख्या सम हो: विषम हाथ मिलाने वाले सम संख्या में आते हैं। (तीन
अतिथियों पर: संभव हस्तमिलन-चित्रों में कभी ठीक एक या तीन विषम प्रविष्टियाँ
नहीं होतीं --- चारों संभव ग्राफ़ जाँच लीजिए।)

\textbf{18.} $n = 1$: $1 = 1$। यदि $1^3 + \dots + n^3 = \left(\frac{n(n+1)}{2}\right)^2$ हो, तो $(n+1)^3$ जोड़ने पर
\[
\frac{n^2(n+1)^2}{4} + (n+1)^3
= \frac{(n+1)^2\left(n^2 + 4n + 4\right)}{4}
= \left(\frac{(n+1)(n+2)}{2}\right)^{\!2} :
\]
वंशानुगति। $n = 3$ के लिए: $1 + 8 + 27 = 36 = 6^2$।

\textbf{19.} $(n, n)$ तक जाने वाला कोई पथ $2n$ क़दम चलता है और
प्रति-विकर्ण $x + y = n$ को ठीक एक जालक-बिंदु $(j, n - j)$ पर पार करता है; उसका पहला आधा
$n$ क़दमों में $j$ बार $R$ वाला पथ है ($\binom nj$ चुनाव), और दूसरा आधा, उल्टा
पढ़ने पर, वैसा ही (सममिति से फिर $\binom nj$)। कटान-बिंदु पर जोड़ने पर: $\binom{2n}{n} = \sum_j \binom nj^2$ ---
वांडरमोंड की सर्वसमिका, खींची हुई।

\textbf{20.} चरण गुणा कीजिए, स्थितियाँ जोड़िए: पट्टिकाएँ और पोकर के हाथ।
चतुराई से लिखिए: पथ $R$--$U$ शब्दों के रूप में, ऑर्डर तारों और छड़ों के रूप
में। दो बार गिनिए: अध्यक्ष-सहित-समितियाँ, हस्तमिलन, बीच से कटे पथ। घटाइए और
सुधारिए: बिखरी टोपियाँ, जिनका अवशेष $\frac1\eu$ है। अगले पड़ाव: \omterm{def:g10:proba:distribution}{प्रायिकता} की भिन्नों
के अधीन यही गिनतियाँ, और दो अध्याय आगे आसन्नता आव्यूहों की पथ-गिनने वाली
घातें।
\end{solution}
